\section{Диагностика согласованных данных UNSW}
\label{app:unsw-diagnostics}
\subsection{Проверенная часть обработки и границы реконструкции}
Восстановленные из публичного зеркала исходные файлы \texttt{UNSW\_NB15\_testing-set.csv} и \texttt{UNSW\_NB15\_training-set.csv} имеют SHA-256, совпадающие с указанными автором. Их объединение в порядке имён файлов, сначала testing, затем training, воспроизводит 257,673 строки исторического кэша. Независимый скрипт повторно вычислил тринадцать числовых кандидатов, два категориальных поля и бинарную метку: все 4,122,768 сравнений значений совпали точно. Во всех исходных строках \texttt{label}=0 соответствует \texttt{attack\_cat}=Normal. Аналогичная проверка TON воспроизвела 16 колонок на 211,043 строках, то есть 3,376,688 значений, без расхождений; метки согласованы с \texttt{type}. Эти проверки не восстанавливают историческое выполнение моделей и не доказывают семантическую эквивалентность полей разных экстракторов. Для новых построчных артефактов идентификатор UNSW должен включать имя файла и номер строки: поле \texttt{id} повторяется между двумя CSV.

Без повторного обучения восстановлены двадцать наборов индексов: два исходных домена для каждого seed от 42 до 51. Они совпали с результатом сохранённой функции балансировки, а размеры и 1,620 сохранённых счётчиков совпали с историческими записями запусков. Каждая объединённая выборка содержит 764 нормальные строки и 3,820 атак. Реконструкция условна на доступный код и среду проверки, включая scikit-learn 1.8.0; исторические хеши выбранных индексов, настроенные преобразователи и совместные модели не сохранились. Поэтому нижеследующие данные являются воспроизведёнными характеристиками выборок, а не восстановленными выходами исторических моделей.

Позднее получены шесть файлов JSONL с 900 записями оценки, включая прежние сетки соотношений классов. Записи согласованы с сохранёнными таблицами отдельных запусков и позволяют независимо пересчитать метрики и агрегаты; это не 900 независимо обученных моделей. Из них 420 записей относятся к текущему ratio5 и подтверждают все 42 текущие суммарные матрицы ошибок. В двенадцати матрицах уже исключённого из выбора ratio50 обнаружен пропуск seed~42; исходные архивные файлы сохранены отдельно от исправленных сумм. Данный дефект прежней сетки не меняет текущие таблицы. Эти контрольные записи не заменяют отсутствующие построчные предсказания и оценки исторических совместных моделей.

\subsection{Категории в сбалансированных источниках и цели}
Для категории $c$ обозначим $P_J=P(attack\mid c)$ в восстановленной совместной выборке, а $P_U$ --- ту же долю в UNSW. Чтобы учесть различие общих долей классов, используем также $LR(c)=P(c\mid attack)/P(c\mid normal)$. Значение $LR>1$ означает более высокую частоту категории внутри класса атак, а $LR<1$ --- внутри нормального класса. Это описательное отношение частот, не параметр классификатора.

\begin{table}[htbp]
\caption{Связи укрупнённых категорий с меткой в восстановленных сбалансированных исходных выборках Joint ($J$) и UNSW ($U$). Для Joint приведены диапазоны по десяти seed; $P=P(attack\mid c)$, $LR=P(c\mid attack)/P(c\mid normal)$. Таблица описывает данные, а не вклад категории в выход обученной модели.}
\label{tab:unsw-category-shift}
\centering\small\setlength{\tabcolsep}{4pt}
\begin{tabular}{lrrrr}
\toprule
Категория & $P_J$ & $P_U$ & $LR_J$ & $LR_U$\\
\midrule
\texttt{proto=other} & 0.0000 & 0.9046 & 0.000 & 5.358\\
\texttt{proto=tcp} & 0.9334--0.9445 & 0.4557 & 2.803--3.406 & 0.473\\
\texttt{proto=udp} & 0.6738--0.6866 & 0.7625 & 0.413--0.438 & 1.813\\
\texttt{state=closed} & 0.8884--0.9105 & 0.4764 & 1.591--2.034 & 0.514\\
\texttt{state=established} & 0.4487--0.4892 & 0.0697 & 0.163--0.192 & 0.042\\
\bottomrule
\end{tabular}
\end{table}

Категория \texttt{other} встречается 35--42 раза в каждой исходной выборке, всегда с меткой normal. В UNSW она содержит 41,916 строк, включая 37,919 атак. Из этих строк 37,760 соответствуют резервному отображению 127 исходных значений протокола, а 4,156 --- явно отображённому ARP. Вся категория и резервное отображение являются разными множествами. Известное модели имя категории не означает одинаковый состав входящих в неё протоколов в разных доменах. В табл.~\ref{tab:unsw-category-shift} направление ассоциации меняется у первых четырёх категорий; для \texttt{established} $LR<1$ в обоих доменах. Поэтому наблюдение в полной TON-популяции о почти исключительно атакующем \texttt{established} нельзя переносить на сбалансированную совместную выборку.

\subsection{Числовые признаки и ограничения простых объяснений}
Девять из тринадцати числовых кандидатов имеют противоположные знаки корреляции Спирмена с меткой в восстановленной смеси и UNSW во всех десяти seed. Семь удовлетворяют также условию $|\rho|>0.05$ в обоих доменах для каждого seed. Например, средняя исходная корреляция duration равна $0.1580$ против $-0.2564$ в UNSW, а flow rate --- $-0.1000$ против $0.2542$. Кандидаты связаны конструктивно и статистически, поэтому эти результаты не являются девятью независимыми механизмами ошибки.

Сохранённые списки выбранных признаков показывают, что \texttt{pkts\_dst} и \texttt{payload\_mean\_dst} не использовались ни в одном из десяти итоговых совместных обучений; \texttt{bytes\_dst} использовался в восьми. Корреляция \texttt{bytes\_dst} в восстановленной смеси слабо отрицательна, как и в UNSW. В общей информативной страте TCP/closed она также отрицательна в обоих полных доменах: $-0.1076$ в TON и $-0.3027$ в UNSW. Следовательно, изменение знака маргинальной корреляции между полными популяциями TON и UNSW может зависеть от состава протоколов и состояний; оно не описывает непосредственно обученную joint-модель. Корреляция также не определяет форму немонотонной функции KAN.

При нулевой длительности формула rate умножает соответствующий счётчик на $10^6$. Однако в восстановленных исходных выборках доля атак при нулевой длительности составляет 77.7--80.2\%, а при положительной --- 84.0--84.4\%; направление относительной связи нулевой длительности с normal сохраняется в UNSW. В UNSW нулевую длительность имеют 3,543 нормальные строки и 64 атаки. Даже исправление всех ошибок исключительно на этих строках могло бы увеличить BA не более чем на $\tfrac12(3543/93000+64/164673)=0.0192427$, то есть на 1.924 процентного пункта. Этого недостаточно для всего наблюдаемого дефицита до BA=0.5. Граница не исключает более сложного влияния rate-признаков на обучение или ошибки на других строках.

Сглаженные отношения $b_s/(p_s+1)$ и $b_d/(p_d+1)$ из~\eqref{eq:harmonized} нельзя использовать как физические средние bytes/packet для проверки MTU. Аналогично, близость или различие распределений duration сами по себе не устанавливают единицы измерения. Связи признаков с меткой подтверждают различия исходных и целевых данных, но для определения их вклада в ошибки нужны исторические настроенные преобразователи, модели и построчные оценки либо отдельно обозначенная новая репликация. В этой редакции нового обучения и настройки адаптации по UNSW не выполнялось.
